%Version 3.1 December 2024
% See section 11 of the User Manual for version history
%
%%%%%%%%%%%%%%%%%%%%%%%%%%%%%%%%%%%%%%%%%%%%%%%%%%%%%%%%%%%%%%%%%%%%%%
%%                                                                 %%
%% Please do not use \input{...} to include other tex files.       %%
%% Submit your LaTeX manuscript as one .tex document.              %%
%%                                                                 %%
%% All additional figures and files should be attached             %%
%% separately and not embedded in the \TeX\ document itself.       %%
%%                                                                 %%
%%%%%%%%%%%%%%%%%%%%%%%%%%%%%%%%%%%%%%%%%%%%%%%%%%%%%%%%%%%%%%%%%%%%%

%%\documentclass[referee,sn-basic]{sn-jnl}% referee option is meant for double line spacing

%%=======================================================%%
%% to print line numbers in the margin use lineno option %%
%%=======================================================%%

%%\documentclass[lineno,pdflatex,sn-basic]{sn-jnl}% Basic Springer Nature Reference Style/Chemistry Reference Style

%%=========================================================================================%%
%% the documentclass is set to pdflatex as default. You can delete it if not appropriate.  %%
%%=========================================================================================%%

%%\documentclass[sn-basic]{sn-jnl}% Basic Springer Nature Reference Style/Chemistry Reference Style

%%Note: the following reference styles support Namedate and Numbered referencing. By default the style follows the most common style. To switch between the options you can add or remove “Numbered” in the optional parenthesis. 
%%The option is available for: sn-basic.bst, sn-chicago.bst%  
 
%%\documentclass[pdflatex,sn-nature]{sn-jnl}% Style for submissions to Nature Portfolio journals
%%\documentclass[pdflatex,sn-basic]{sn-jnl}% Basic Springer Nature Reference Style/Chemistry Reference Style
\documentclass[pdflatex,sn-mathphys-num]{sn-jnl}% Math and Physical Sciences Numbered Reference Style
%%\documentclass[pdflatex,sn-mathphys-ay]{sn-jnl}% Math and Physical Sciences Author Year Reference Style
%%\documentclass[pdflatex,sn-aps]{sn-jnl}% American Physical Society (APS) Reference Style
%%\documentclass[pdflatex,sn-vancouver-num]{sn-jnl}% Vancouver Numbered Reference Style
%%\documentclass[pdflatex,sn-vancouver-ay]{sn-jnl}% Vancouver Author Year Reference Style
%%\documentclass[pdflatex,sn-apa]{sn-jnl}% APA Reference Style
%%\documentclass[pdflatex,sn-chicago]{sn-jnl}% Chicago-based Humanities Reference Style

%%%% Standard Packages
%%<additional latex packages if required can be included here>

\usepackage{graphicx}%
\usepackage{multirow}%
\usepackage{amsmath,amssymb,amsfonts}%
\usepackage{amsthm}%
\usepackage{mathrsfs}%
\usepackage[title]{appendix}%
\usepackage{xcolor}%
\usepackage{textcomp}%
\usepackage{manyfoot}%
\usepackage{booktabs}%
\usepackage{algorithm}%
\usepackage{algorithmicx}%
\usepackage{algpseudocode}%
\usepackage{listings}%
%%%%

%%%%%=============================================================================%%%%
%%%%  Remarks: This template is provided to aid authors with the preparation
%%%%  of original research articles intended for submission to journals published 
%%%%  by Springer Nature. The guidance has been prepared in partnership with 
%%%%  production teams to conform to Springer Nature technical requirements. 
%%%%  Editorial and presentation requirements differ among journal portfolios and 
%%%%  research disciplines. You may find sections in this template are irrelevant 
%%%%  to your work and are empowered to omit any such section if allowed by the 
%%%%  journal you intend to submit to. The submission guidelines and policies 
%%%%  of the journal take precedence. A detailed User Manual is available in the 
%%%%  template package for technical guidance.
%%%%%=============================================================================%%%%

%% as per the requirement new theorem styles can be included as shown below
\theoremstyle{thmstyleone}%
\newtheorem{theorem}{Theorem}%  meant for continuous numbers
%%\newtheorem{theorem}{Theorem}[section]% meant for sectionwise numbers
%% optional argument [theorem] produces theorem numbering sequence instead of independent numbers for Proposition
\newtheorem{proposition}[theorem]{Proposition}% 
%%\newtheorem{proposition}{Proposition}% to get separate numbers for theorem and proposition etc.

\theoremstyle{thmstyletwo}%
\newtheorem{example}{Example}%
\newtheorem{remark}{Remark}%

\theoremstyle{thmstylethree}%
\newtheorem{definition}{Definition}%

\raggedbottom
%%\unnumbered% uncomment this for unnumbered level heads

\begin{document}

\title[HAG-DTA]{HAG-DTA: \ drug-target binding affinity prediction with hierarchical multi-scale attention fusion graph neural networks}

%%=============================================================%%
%% GivenName	-> \fnm{Joergen W.}
%% Particle	-> \spfx{van der} -> surname prefix
%% FamilyName	-> \sur{Ploeg}
%% Suffix	-> \sfx{IV}
%% \author*[1,2]{\fnm{Joergen W.} \spfx{van der} \sur{Ploeg} 
%%  \sfx{IV}}\email{iauthor@gmail.com}
%%=============================================================%%

\author[1]{\fnm{Jiannuo} \sur{Li}}
\author[1]{\fnm{Longhao} \sur{Lei}}
\author*[1]{\fnm{Lan} \sur{Yao}}\email{yao@hnu.edu.cn}

\affil[1]{\orgdiv{School of Mathematics}, \orgname{University of Hunan}, \orgaddress{ \city{Changsha}, \postcode{410082}, \country{China}}}

%%==================================%%
%% Sample for unstructured abstract %%
%%==================================%%

\abstract{\textbf{Motivation}: Accurate prediction of the affinity between drugs and target proteins is important for the drug development process, and deep learning methods have emerged as significant tools in drug-target affinity prediction (DTA). However, the existing DTA models ignore the impact of specific structural features within drug molecules, leading to inadequate representation of the drug molecule characteristics.\newline
\textbf{Results}: We propose the HAG-DTA model to improve the DTA prediction accuracy by sufficiently utilizing multi-scale features of drugs, in which the multi-scale features are extracted by two modules. One module is the hierarchical graph neural network (HGNN) for learning the local features, including the inter and intra structures of functional groups. Another is deep graph neural network (DGNN) for capturing the global features of drugs. Meanwhile, we introduced MMD loss to limit the distribution difference between local and global features. Furthermore, a convolutional neural network (CNN) with multiple convolutional filter sizes is introduced to extract protein features. Finally, attention mechanisms are employed to effectively integrate the drug features with the target protein features. The experimental results on four benchmark datasets show that the HAG-DTA model outperforms baseline models on DTA prediction and interpretability.}

%%================================%%
%% Sample for structured abstract %%
%%================================%%

% \abstract{\textbf{Purpose:} The abstract serves both as a general introduction to the topic and as a brief, non-technical summary of the main results and their implications. The abstract must not include subheadings (unless expressly permitted in the journal's Instructions to Authors), equations or citations. As a guide the abstract should not exceed 200 words. Most journals do not set a hard limit however authors are advised to check the author instructions for the journal they are submitting to.
% 
% \textbf{Methods:} The abstract serves both as a general introduction to the topic and as a brief, non-technical summary of the main results and their implications. The abstract must not include subheadings (unless expressly permitted in the journal's Instructions to Authors), equations or citations. As a guide the abstract should not exceed 200 words. Most journals do not set a hard limit however authors are advised to check the author instructions for the journal they are submitting to.
% 
% \textbf{Results:} The abstract serves both as a general introduction to the topic and as a brief, non-technical summary of the main results and their implications. The abstract must not include subheadings (unless expressly permitted in the journal's Instructions to Authors), equations or citations. As a guide the abstract should not exceed 200 words. Most journals do not set a hard limit however authors are advised to check the author instructions for the journal they are submitting to.
% 
% \textbf{Conclusion:} The abstract serves both as a general introduction to the topic and as a brief, non-technical summary of the main results and their implications. The abstract must not include subheadings (unless expressly permitted in the journal's Instructions to Authors), equations or citations. As a guide the abstract should not exceed 200 words. Most journals do not set a hard limit however authors are advised to check the author instructions for the journal they are submitting to.}

\keywords{Drug-target Binding Affinity, Multi-scale feature, Hierarchical Network, Graph neural network}

%%\pacs[JEL Classification]{D8, H51}

%%\pacs[MSC Classification]{35A01, 65L10, 65L12, 65L20, 65L70}

\maketitle

\section{Introduction}

Drug-target affinity (DTA) prediction plays a critical role in computer-aided drug discovery, as it enables the estimation of binding strength between small-molecule compounds and target proteins. Accurate DTA prediction is essential for a wide range of downstream applications, including virtual screening, drug repositioning, and lead compound optimization. However, experimental determination of binding affinity is often time-consuming, labor-intensive, and costly \cite{ozturk2018deepdta}. In this context, computational approaches for DTA prediction have become increasingly important, as they allow large-scale screening of candidate compounds and help prioritize promising drug-target pairs for further experimental validation.

With the rapid development of deep learning, various neural network-based methods have been proposed to improve DTA prediction. Early deep learning models encode drugs and proteins as one-dimensional sequences, enabling end-to-end representation learning, such as DeepDTA \cite{ozturk2018deepdta}. More recently, graph neural network (GNN)-based approaches have been widely adopted to represent drugs as molecular graphs, such as GraphDTA \cite{nguyen2021graphdta}, allowing the incorporation of structural information into DTA prediction. These methods have demonstrated strong performance improvements by capturing local topological patterns of molecules and sequence-level features of proteins.

Despite these advances, several limitations remain. First, most existing methods focus on single-scale molecular representations or implicitly capture multi-scale information without explicitly modeling the hierarchical structure of molecules. In practice, drug molecules naturally exhibit hierarchical organization, where atoms form functional groups and functional groups further compose molecular fragments. Although some studies have explored multi-scale graph learning \cite{yang2022mgraphdta,ying2018hierarchical}, the explicit modeling of hierarchical molecular relationships remains insufficient, which may limit the ability of models to capture chemically meaningful patterns associated with drug-target binding.

Second, the fusion of drug and protein features is often implemented through simple concatenation or loosely coupled attention mechanisms \cite{zhao2023attentiondta}. Such strategies may not fully exploit the complementary nature of molecular structures and protein sequences, and they often lack an effective mechanism to dynamically assess the importance of different feature sources. As a result, the interaction between drug and protein representations may not be sufficiently captured.

Third, interpretability remains a critical yet underexplored aspect in many DTA prediction frameworks. In practical drug discovery, the goal is not only to predict binding affinity accurately, but also to identify key atoms, functional groups, or molecular fragments that contribute to binding. However, existing deep learning-based methods often act as black-box models, and their decision-making processes are difficult to interpret. Although attention-based mechanisms have been introduced to provide partial explanations \cite{zhao2023attentiondta}, they typically offer coarse-grained insights and do not explicitly reveal how molecular information is organized across different structural levels.

To address these challenges, we propose a Hierarchical Attention-based Graph model for Drug-Target Affinity prediction (HAG-DTA). The proposed model explicitly learns hierarchical molecular representations by organizing drug features into atom-level, functional group-level, and fragment-level structures. In addition, a deep graph representation branch is introduced to capture global molecular topology, while a convolutional neural network is employed to extract protein sequence features. These heterogeneous features are then integrated through an attention-based fusion mechanism, which dynamically assigns importance to different feature components. By combining hierarchical structural modeling with adaptive feature fusion, the proposed model not only improves prediction performance but also enhances structural interpretability.

Extensive experiments on multiple benchmark datasets demonstrate that HAG-DTA achieves competitive performance in both regression and classification tasks. Furthermore, ablation studies, hyperparameter analyses, and visualization results provide additional insights into the effectiveness and interpretability of the proposed model.

\section{Related Work}\label{sec2}

\subsection{Drug-Target Affinity Prediction Methods}

Drug-target affinity (DTA) prediction is a fundamental task in computer-aided drug design. It aims to estimate the binding strength between small-molecule compounds and target proteins, thereby supporting virtual screening, drug repositioning, and candidate compound optimization. Early DTA prediction studies mainly relied on handcrafted molecular descriptors, protein sequence similarities, and traditional machine learning models. Representative methods include KronRLS \cite{pahikkala2015toward}, which formulates drug-target interaction prediction with Kronecker regularized least squares, and SimBoost \cite{he2017simboost}, which predicts continuous binding affinity values using gradient boosting with drug-target similarity features. However, the performance of these methods largely depends on feature engineering, and their ability to capture complex drug-target interactions is limited.

With the rapid development of deep learning, neural networks have been increasingly used to automatically learn expressive representations of drugs and proteins. DeepDTA \cite{ozturk2018deepdta} first applied convolutional neural networks to DTA prediction by encoding SMILES-based drug strings and protein sequences in an end-to-end manner. WideDTA \cite{ozturk2019widedta} further introduced multiple textual representations of drugs and proteins, while AttentionDTA \cite{zhao2023attentiondta} incorporated attention mechanisms to enhance the modeling of important sequence regions. In addition, TransformerCPI \cite{chen2020transformercpi} and MolTrans \cite{huang2021moltrans} explored self-attention and Transformer-based architectures to capture compound-protein interaction patterns. These methods demonstrate the effectiveness of deep learning in DTA prediction and reduce the dependence on manually designed features.

Nevertheless, sequence-based methods usually treat SMILES strings as one-dimensional sequences. As a result, they cannot fully exploit the intrinsic topological structures of molecules. To address this limitation, graph neural network (GNN)-based models have been widely adopted in recent years. GraphDTA \cite{nguyen2021graphdta} represents drugs as molecular graphs and uses graph convolutional networks to learn drug features, achieving improved prediction performance. Based on this idea, MGraphDTA and DGraphDTA \cite{yang2022mgraphdta,jiang2022sequence} further incorporate multi-scale topological information and spatial structural features to improve molecular representation. AttentionMGT-DTA \cite{wu2024attentionmgtdta} further combines graph Transformer and attention mechanisms for multimodal DTA prediction. These studies indicate that molecular graph modeling is essential for accurate DTA prediction.

Recently, more advanced DTA models have been proposed to further improve molecular and protein representation learning. GPCNDTA \cite{zhang2023gpcndta} introduces pharmacophore prior information and modal interaction strategies to enhance drug-target feature extraction. GEFormerDTA \cite{liu2024geformerdta} introduces a Transformer graph-based early fusion strategy and reports competitive results on the Davis and KIBA benchmark datasets. GNPDTA \cite{ye2024graphpre} employs graph neural pre-training to extract low-level features from drug atom graphs and target residue graphs. GLCN-DTA \cite{qi2024extendedgraph} enhances DTA prediction by combining graph learning operations with graph convolution operations. CSCo-DTA \cite{wang2024crossscale} integrates molecule-scale and network-scale information through cross-scale graph contrastive learning. MDCT-DTA \cite{zhu2024multiscale} introduces multi-scale diffusion and interactive learning for DTA prediction. WPGraphDTA \cite{hu2025powergraph} constructs drug-target affinity representations based on power graph and Word2Vec. MEGDTA \cite{hou2025megdta} further incorporates protein three-dimensional structural information and ensemble graph neural networks to capture diverse drug and target features. These recent studies indicate that richer structural information, pre-training strategies, contrastive learning, and more effective fusion mechanisms have become important directions in DTA prediction. However, how to explicitly organize molecular representations across atom-level, functional group-level, and fragment-level structures remains insufficiently explored.

\subsection{Graph Neural Network-Based Molecular Representation Learning}

Graph neural networks have become powerful tools for molecular representation learning. They can directly model atoms as nodes and chemical bonds as edges. In this way, both atom attributes and molecular topology can be captured.

Graph Convolutional Networks (GCNs) \cite{kipf2016semi} update node representations through normalized neighborhood aggregation. Each node aggregates information from itself and its neighboring nodes, which enables the model to capture local structural patterns. GCNs have a simple architecture and stable training behavior. However, when many layers are stacked, node representations may become overly similar, leading to the over-smoothing problem.

Graph Isomorphism Networks (GINs) \cite{xu2018powerful} improve the discriminative power of GNNs by using sum-based aggregation and multilayer perceptrons. This design allows GINs to preserve neighborhood multiset information and provides strong expressive capability comparable to the Weisfeiler-Lehman graph isomorphism test. Therefore, GINs are suitable for molecular tasks that require fine-grained structural distinction. However, they may also be sensitive to noise and prone to overfitting.

Graph Attention Networks (GATs) \cite{velivckovic2017graph} introduce an attention mechanism into neighborhood aggregation. By assigning different attention weights to neighboring nodes, GATs can adaptively focus on more informative local structures. This improves modeling flexibility and provides a certain degree of interpretability. The main limitation is that GATs usually require higher computational costs.

GraphSAGE \cite{hamilton2017inductive} learns inductive node representations through neighborhood sampling and aggregation. It does not require access to the entire graph during training. This property enables GraphSAGE to generalize to unseen nodes or graphs, making it useful for large-scale and dynamic graph scenarios. Its performance is closely related to the sampling strategy and aggregation function.

To further enhance molecular representation learning, several studies have explored multi-scale and hierarchical graph modeling strategies \cite{ying2018hierarchical,yang2022mgraphdta,wang2024crossscale,qi2024extendedgraph}. These methods aim to capture both local substructures and global topological information. DiffPool \cite{ying2018hierarchical} provides a differentiable graph pooling framework that learns node-to-cluster assignments and generates hierarchical graph representations. MGraphDTA \cite{yang2022mgraphdta} employs deep multiscale graph neural networks to learn molecular representations from different receptive fields. CSCo-DTA \cite{wang2024crossscale} integrates molecule-scale and network-scale information through cross-scale graph contrastive learning, while GLCN-DTA \cite{qi2024extendedgraph} improves DTA prediction by extending graph learning operations. These studies suggest that multi-scale graph modeling is important for capturing complex molecular structures.

Despite these advances, most existing DTA models still focus on single-scale graph representations or shallow hierarchical structures. Although multi-scale graph learning can expand the receptive field of molecular representation, the fine-grained relationships among atoms, functional groups, and molecular fragments are not sufficiently modeled. In particular, many methods lack an explicit mechanism to learn how atom-level information is progressively organized into higher-level molecular substructures. This limitation may reduce the ability of existing models to represent chemically meaningful structural patterns related to drug-target binding.

\subsection{Multimodal Fusion and Model Interpretability Analysis}

In DTA prediction, drugs and proteins belong to heterogeneous modalities. Effective feature fusion is therefore essential for modeling drug-target interactions. Existing studies commonly fuse drug and protein representations through feature concatenation or attention mechanisms. For example, AttentionDTA \cite{zhao2023attentiondta} uses attention weights to highlight important features and improve prediction accuracy. TransformerCPI \cite{chen2020transformercpi} applies self-attention mechanisms to learn compound-protein interaction features, while MolTrans \cite{huang2021moltrans} introduces a knowledge-inspired substructure mining strategy and interaction modeling module. AttentionMGT-DTA \cite{wu2024attentionmgtdta} further adopts graph Transformer and attention mechanisms to integrate multimodal drug and target information. These studies show that attention-based fusion can improve the interaction modeling ability of DTA models. However, simple concatenation may not sufficiently capture the relative contribution of different drug and protein features, while attention-based fusion still requires more explicit structural guidance.

In addition, multimodal learning has also been applied to other pharmaceutical prediction tasks. DPSP \cite{masumshah2023dpsp} and NNPS \cite{masumshah2021polypharmacyside}, for example, propose deep learning frameworks for polypharmacy side-effect prediction, showing the potential of integrating heterogeneous drug-related information. These studies indicate that the combination of chemical structures, biological targets, side effects, and pathway information can improve the representation ability of deep learning models in drug-related prediction tasks.

Model interpretability has also attracted increasing attention. In DTA prediction, interpretability is important because it can help identify key atoms, functional groups, or protein regions that contribute to binding affinity. MONN \cite{li2020monn} predicts both compound-protein interactions and binding affinities, and explores attention-based interpretation of molecular interactions. TransformerCPI \cite{chen2020transformercpi} maps attention weights to protein sequences and compound atoms to support interaction interpretation. MolTrans \cite{huang2021moltrans} further provides interpretable substructure-level interaction modeling. Attention weight analysis and molecular substructure visualization are commonly used for this purpose. These techniques can provide useful clues for rational drug design.

Although existing methods have improved feature fusion and interpretability, several limitations remain. First, multi-scale molecular structural information is still insufficiently utilized. Second, the correlations among different structural levels are rarely modeled explicitly. Third, the interpretation of model decision-making is often coarse and lacks detailed structural evidence. These limitations motivate the development of models that can jointly learn hierarchical molecular representations, integrate drug and protein features adaptively, and provide more informative structural visualization.

To address these challenges, we propose a Hierarchical Attention-based Graph model for Drug-Target Affinity prediction (HAG-DTA). In addition to improving prediction accuracy, HAG-DTA is designed to enhance structural interpretability by explicitly modeling hierarchical molecular representations. Specifically, the model progressively learns atom-level, functional group-level, and molecular fragment-level features, allowing the contribution of different molecular substructures to be visualized and analyzed. These hierarchical drug representations are further integrated with global molecular topology and protein sequence features through an attention-based fusion mechanism.


\section{Materials and methods}\label{sec3}

The purpose of this study is to extract the multiscale features of drug molecules to fuse molecular and target features to predict DTA(Drug-Target Affinity). In light of this, we propose HAG-DTA, an end-to-end prediction algorithm. HAG-DTA takes the SMILES notation of drug molecules and protein sequences as inputs, and yields affinity values as outputs, which are determined by dissociation constants or KIBA (kinase inhibitor bioactivity data) scores. The HAG-DTA model consists of two modules: the hierarchical GNN network and the deep GNN network, as illustrated in \autoref{fig1} The hierarchical GNN network includes two modules, namely the 
$H_1$ Block and $H_2$ Block, each comprising an encoding block and a pooling block. The $H_1$ Block is responsible for extracting the features of specific functional groups $x_{H_1}$, while the $H_2$ Block captures the characteristics of molecular fragments $x_{H_2}$. The deep GNN network is composed of five GNN layers and utilizes gated skip connections to extract the global features $x_G$ of drug molecules. Furthermore, we construct a Convolutional Neural Network (CNN) module to extract protein features $x_P$. This module comprises three parallel one-dimensional convolutional neural networks and a max pooling layer. During the prediction phase, the HAG-DTA model combines the multiscale features of molecules and target features using an attention mechanism and feeds them into an MLP layer to obtain prediction values.

\begin{figure*}[!t]
\centering
\includegraphics[width=\linewidth]{image/figure1.pdf}
\caption{The figure shows the HAG-DTA architecture. In the figure, the Hierarchical Network (HGNN), Deep Graph Network (DGNN), and Convolutional Neural Network (CNN) are employed to extract multi-scale features from drug molecules and proteins. The Attention Block is utilized to fuse the extracted multi-scale features. Finally, the fused features are passed through several fully connected layers to estimate the drug-target affinity values.}
\label{fig1}
\end{figure*}
\subsection{Input representation}\label{subsubsec1}
The drug molecules were originally stored in the SMILES(Simplified Molecular Input Line Entry System) format, from which the descriptors of the drug can be inferred. In this study, we adopt a method consistent with GraphDTA\cite{nguyen2021graphdta} to convert drug molecules into undirected graphs. Specifically, we use the RDKit\cite{landrum2006rdkit} tool to treat atoms in the drug molecule as nodes and chemical bonds as edges. Consequently, the drug molecule is transformed into an undirected graph, from which the adjacency matrix A and original features X are extracted. \autoref{fig2}(b) illustrates the transformation process. 
\newline
\hspace*{1em}The input target represents a protein sequence, with each letter corresponding to an amino acid. We establish a mapping between each letter and a specific integer, effectively converting the protein sequence into an integer sequence. For the sake of facilitating model comparison, our mapping method in this study aligns with GraphDTA\cite{nguyen2021graphdta}, and the maximum length of the protein sequence is set to 1200.
\subsection{Multi-scale feature fusion}
The hierarchical GNN network utilizes hierarchical clustering to map the nodes of molecular graphs to a set of clusters. Subsequently, these clusters are reconstructed as nodes into a coarsened graph that encompasses functional group features or molecular fragment features. \autoref{fig2}(c) illustrates the mapping process employed by the hierarchical GNN network. The diagram depicts the $H_1$ Block clustering closely connected atoms within the input drug molecule and reconstructing them as a coarsened graph $G_1$. The nodes within this $G_1$ graph incorporate the features of functional groups. The $H_2$ Block takes $G_1$ as input and captures the connection relationships between functional groups to generate the coarsened graph $G_2$. The nodes in $G_2$ contain fragment features.
\newline
\hspace*{1em}\autoref{fig2}(a) provides a visualization of how deep GNN networks extract global features from drug molecules. To accomplish this, deep GNN networks employ gated skip connections, enabling the fusion of multi-order information and the extraction of global features. 
\newline
\hspace*{1em}\autoref{fig2}(d) visualizes the process of extracting protein features using the Convolutional Neural Network (CNN) module.
\newline
\hspace*{1em}The utilization of attention mechanisms offers the ability to dynamically adjust the significance of features, thereby enhancing the model's generalization performance. Hence, after obtaining the multi-scale features $x_G$, $x_{H_1}$ and $x_{H_2}$ of the drug, as well as the protein feature $x_p$, we employ an attention mechanism for feature fusion. Firstly, we construct two linear layers to calculate the weight coefficients of the four features: 
$$w_{j}=\operatorname{Linear}_{j}\left(x_{j}\right)=\sigma\left(x_{j} W_{1, j}+b_{1, j}\right) W_{2, j}, j \in\left\{G, H_{1}, H_{2}, P\right\},$$
where $\sigma\left(\cdot \right)$ represents the activation function, $W_{1, j}$, $W_{2, j}$, and $b_{1, j}$ are the learnable parameters in the linear layer, and $w_G$, $w_{H_1}$, $w_{H_2}$, and $w_p$ are the weight coefficients of features $x_G$, $x_{H_1}$, $x_{H_2}$ and $x_p$, respectively.
\newline
\hspace*{1em}Then we use the softmax function to normalize the weight coefficients, obtaining attention scores $\alpha_G$, $\alpha_{H_1}$, $\alpha_{H_2}$ and $\alpha_p$: 
$$
\alpha_{j}=\frac{e^{w_{j}}}{e^{w_{G}}+e^{w_{H_{1}}}+e^{w_{H_{2}}}+e^{w_{P}}} \quad j \in\left\{G, H_{1}, H_{2}, P\right\}
.$$
\hspace*{1em}Finally, based on these attention scores, we fuse the features and input the fused features into the MLP for prediction, resulting in an affinity value denoted by:
$$
\mathrm{y}=\operatorname{MLP}\left(\alpha_{G} x_{G}\left\|\alpha_{H_{1}} x_{H_{1}}\right\| \alpha_{H_{2}} x_{H_{2}} \| \alpha_{P} x_{P}\right)
,$$
where $\|$ signifies the concatenate operation.
\begin{figure}[!t]%
\centering
\includegraphics[width=\linewidth]{image/figure2.pdf}
\caption{Figure (b) illustrates the process of converting drug SMILES encoding into molecular graphs using RDKit. Figure (c) visualizes the process of extracting local features from drug molecules using HGNN,  where the $G_1$ and $G_2$ contain functional groups and molecular fragment features, respectively.}\label{fig2}
\end{figure}
\subsection{Graph neural network layer}
Let $G=\{V,E\}$ be a graph of a given drug molecule, where $V$ is the set of nodes, $E$ is the set of edges, $v_i \in V$ is the i-th atom, and $e_{ij}\in E$ is the chemical bond between the i-th and j-th atoms. The node set $V$ is represented by a node feature matrix $X\in \mathbb{R}^{n\times d}$, and the edge set $E$ is represented by an adjacency matrix $A\in \mathbb{R}^{n\times n}$, where $n$ is the number of atoms, $d$ is the dimensionality of node features.
\newline
\hspace*{1em}The HAG-DTA model extensively employs GNN layers, which serve as a neighbor aggregation strategy to iteratively aggregate and transfer information about connected nodes in the graph. GNN typically consists of two stages: message passing and readout. Different design approaches give rise to various GNN variants. In this study, we will evaluate four representative GNN variants: GCN\cite{kipf2016semi}, SAGE\cite{hamilton2017inductive}, GAT\cite{velivckovic2017graph} and GIN\cite{xu2018powerful}.
\newline
\hspace*{1em}The GCN model is defined as:
$$x_{i}^{(l)}=\sigma\left(\sum_{j \in N(i) \cup\{i\}} W^{(l)} \frac{x_{i}^{(l-1)}}{\sqrt{|N(i)||N(j)|}}\right),$$
where $W^{(l)} \in \mathbb{R}^{h\times d}$ is the learnable weight matrix, $N(i)$ represents the set of neighbors of node $i$, $\sigma\left( \cdot \right)$ denotes the ReLU activation function, and $x_{i}^{(l-1)}$ represents the features of the i-th node in the output of the $l-1$ layer.
\newline
\hspace*{1em}The SAGE model is defined as:
$$
x_i^{(l)}=\sigma(W^{(l)}\cdot CONCAT(x_i^{(l-1)},AGG(\{x_j^{(l-1)},\forall j\in N(i)\}))
,$$
SAGE\cite{hamilton2017inductive} provides three functions: Mean aggregator, LSTM aggregator, and Pooling aggregator. In this study, the Mean aggregator will be used.
\newline
\hspace*{1em}The GAT model is defined as:
$$
x_i^{(l)}=\sigma(\alpha_{i,i}W^{(l)}x_i^{(l-1)}+\sum_{j\in N(i)}\alpha_{i,j}W^{(l)}x_j^{(l-1)}),
$$
where the calculation method for the attention coefficient $\alpha_{i,j}$ is defined as:
$$
\alpha_{i,j}=\frac{\exp(LeakyReLU(a^T[W^{(l)}x_i^{(l-1)}||W^{(l)}x_j^{(l-1)}]))}{\Sigma_{k\in N(i)\cup\{i\}}\exp(LeakyReLU(a^T[W^{(l)}x_i^{(l-1)}||W^{(l)}x_k^{(l-1)}]))}
,$$
Here, matrix $a\in\mathbb{R}^{2h\times 1}$ maps the attention weights to $\mathbb{R}$.
\newline
\hspace*{1em}The GIN model is defined as:
$$x_i^{(l)}=MLP((1+\epsilon^{(l)})\cdotp x_i^{(l-1)}+\sum_{j \in N(i)}x_j^{(l-1)}),$$
where MLP represents fully-connected neural network layers, $\epsilon^{(l)}$ is either a learnable parameter of layer $l$ or a fixed scalar.
\subsection{Hierarchical pooling}
In the hierarchical GNN network in HAG-DTA, we introduce a hierarchical clustering strategy to extract substructure features of drug molecules. The hierarchical GNN network consists of encoding blocks and pooling blocks. The structure of the encoding blocks and pooling blocks is similar, composed of GNN layers and GraphNorm layers, but they serve different purposes.
\newline
\hspace*{1em}The encoding block encodes the node feature matrix $x^{(l)}$ and outputs an embedding matrix $Z^{(l)}\in\mathbb{R}^{\mathrm{n_l\times d_l}}$, where $n_l$ represents the number of input nodes in the lth layer, and $d_l$ represents the dimension of the output nodes in the lth layer. The embedding matrix $Z^{(l)}$ can be calculated as: 
$$
Z^{(l)}=GNN_{l,embed}(A^{(l)},X^{(l)})
.$$
\hspace*{1em}The pooling block outputs a cluster allocation matrix $S^{(l)}\in\mathbb{R}^{\mathrm{n_l\times n_{l+1}}}$, where $n_{l+1}$ is the pre-defined number of clusters for the output in the lth layer and serves as a hyperparameter of the model.
\begin{equation}
S^{(l)}=softmax(GNN_{l,pool}(A^{(l)},X^{(l)})),\label{equation1}
\end{equation}
where the ith row of $S^{(l)}$ represents the probability of the lth layer nodes being assigned to the ith clusters. In equation \ref{equation1}, $A^{(l)}$ and $X^{(l)}$ respectively represent the adjacency matrix and node features of the coarsened graph in the lth layer.
\newline
\hspace*{1em}The cluster allocation matrix $S^{(l)}$, obtained along with the embedding matrix $Z^{(l)}$, is subsequently fed into the DiffPool layer for the purpose of coarsening the input graph.
$$
(A^{(l+1)},X^{(l+1)})=\mathrm{DiffPool}(A^{(l)},Z^{(l)},S^{(l)}).
$$
\hspace*{1em}In the DiffPool layer, the specific calculation method is defined as equation:
$$
X^{(l+1)}=S^{(l)}Z^{(l)}\in\mathbb{R}^{\mathrm{n}_{l+1}\times\mathrm{d}},
$$
$$
A^{(l+1)}=S^{(l)}A^{(l)}S^{(l)}\in \mathbb{R} ^{n_{l+1}\times n_{l+1}},
$$
where $X^{(l+1)}$ and $A^{(l+1)}$ respectively represent the feature matrix and adjacency matrix of the coarsened graph.
\subsection{Deep graph neural network}
To effectively aggregate information from distant neighbors in a given central atom, it is necessary to stack multiple GNN layers. However, constructing a deep GNN network can result in over-smoothing and the vanishing gradient problem\cite{li2019deepgcns}. Consequently, most state-of-the-art GNN models do not exceed three layers. Nevertheless, shallow GNN networks fail to capture the global information of molecules. To overcome this limitation, we introduce a gated skip connection mechanism to construct a deep GNN network. This deep GNN network comprises five GNN layers, leveraging both the message passing mechanism\cite{ryu2018deeply} and the gated skip-connection mechanism to extract global features by integrating multi-order information from molecules, as depicted in \autoref{fig2}(c). The gated skip connection mechanism enables the model to adaptively blend features of varying scales, thus mitigating the issues of over-smoothing and the vanishing gradient. The gated skip-connection mechanism is defined as:
$$x_i^{(l+1)}=z_i\odot x_i^{(l+1)}+(1-z_i)\odot x_i^{(l)},$$
$$
z_i=\sigma(W_{z,1}x_i^{(l+1)}+W_{z,2}x_i^{(l)}+b_z),$$
where $\odot$ denotes the Hadamard product, $W_{z,1}$, $W_{z,2}$ and $b_z$ are trainable parameters, $\sigma(\cdot )$ denotes the activation function, and $z_i$ denotes the information retention ratio coefficient of the previous GNN layer.
\subsection{Loss function}
For the affinity prediction task, we utilize mean squared error (MSE) as the loss function:
$$\mathcal{L}_{MSE}=\frac1n\sum_{i=1}^n(\tilde{y}_i-y_i)^2,$$
where $p$ refers to the prediction vector, while $y$ represents the actual output vector, and $n$ is the number of samples.
\newline
\hspace*{1em}To prevent the model from getting trapped in local minima during the training of the hierarchical GNN network, apart from using the Mean Squared Error (MSE) loss function, we will incorporate auxiliary link prediction objectives which is defined as equation \ref{equation2} and entropy regularization which is defined as equation \ref{equation3}:
\begin{equation}
\mathcal{L}_{LP}=||A^{(l)}-S^{(l)}S^{(l)T}||_F,\label{equation2}
\end{equation}
\begin{equation}
\mathcal{L}_E=\frac1n\sum_{i=1}^nH(S_i^{(l)}),\label{equation3}
\end{equation}
where $||\cdot||_F$ denotes the Frobenius paradigm and $H(\cdot)$ denotes the entropy function.
\newline
\hspace*{1em}The differences in the extraction models for the local feature $x_{H_2}$ and the global feature $x_G$ are mainly in the pooling, and thus can be considered as small differences in the feature distributions of these two features. Therefore, in order to accelerate the convergence of the model and improve the prediction efficiency, we introduce the MMD\cite{gretton2012kernel}(Maximum Mean Discrepancy) loss function, which is defined as  equation \ref{equation4}:
\begin{equation}
\begin{split}
\mathcal{L}_{MMD} = & \frac{1}{n(n-1)} \sum_{i=1}^{n} \sum_{j \neq i}^{n} k\left(x^{(i)}_G, x^{(j)}_G\right)  + \frac{1}{n(n-1)} \sum_{i=1}^{n} \sum_{j \neq i}^{n} k\left(x^{(i)}_{H_2}, x^{(j)}_{H_2}\right) \\
& - \frac{2}{n^2} \sum_{i=1}^{n} \sum_{j=1}^{n} k\left(x^{(i)}_G, x^{(j)}_{H_2}\right), \label{equation4}
\end{split}
\end{equation}
where $x^{(i)}_G$ and $x^{(i)}_{H_2}$ are the global and local features of the $i$th drug molecule, and $k(\cdot,\cdot)$ is the Gaussian kernel function.
\newline
\hspace*{1em}Therefore, combining the four formulas, the loss function for HAG-DTA model training is defined as:
$$
\mathcal{L}_{loss}=\mathcal{L}_{MSE}+\alpha(\mathcal{L}_{LP}+\mathcal{L}_E) + \beta \mathcal{L}_{MMD}
.$$

It is important to note that in order to reduce the difference between two distributions while retaining certain differences, we will set the regularization coefficient $\beta$ of the $\mathcal{L}_{MMD}$ term to a small value.
\section{Results and Discussion}
\subsection{Datasets and metrics}
In this study, we conducted a comprehensive performance evaluation of the HAG-DTA model on four widely recognized and publicly available datasets: Davi\cite{davis2011comprehensive}, Kiba\cite{tang2014making}, Human\cite{liu2015improving}, and \textit{C.elegans}\cite{liu2015improving}.
\newline
\hspace*{1em}The Davis dataset comprises 442 kinase proteins and their corresponding 68 inhibitors, with binding affinities measured by the dissociation constant (KD) ranging from 0.016 to 10,000. Considering the wide range of values, we used equation \ref{equation6} to convert them into logarithmic values. The Kiba dataset includes binding affinities of 2,116 drugs and 229 targets, measured by the KIBA score, ranging from 0.0 to 17.2. In terms of the number of drug-target pairs with interactions, the Kiba dataset is approximately four times larger than the Davis dataset.
\begin{equation}
\mathrm{pK_d=-\log_{10}(\frac{K_d}{10^9})}.\label{equation6}
\end{equation}
\hspace*{1em}The Human and \textit{C.elegans} datasets consist of high-confidence negative samples and positive samples retrieved from the Drugbank database and the Matador dataset. \autoref{table1} provides more detailed information about these datasets.
\begin{table}[!t]
\caption{Summary of the four datasets\label{table1}}%
\begin{tabular*}{\columnwidth}{@{\extracolsep\fill}lllll@{\extracolsep\fill}}
\toprule
Dataset   & Task types     & Compounds & Proteins & Interactions \\ 
\midrule
Davis     & Regression     & 68        & 442      & 30056 \\
Kiba      & Regression     & 2111      & 229      & 118254    \\
Human     & Classification & 2726      & 2001     & 6728     \\
\textit{C.elegans} & Classification & 1767      & 1876     & 7785   \\
\botrule
\end{tabular*}
\end{table}
\newline
\hspace*{1em}To compare the performance against the baseline model in the regression task, we evaluated the model using three metrics: Mean Squared Error (MSE), $r_m^2$ and Concordance Index (CI)\cite{gonen2005concordance}. The CI index measures the discriminative ability among different models on a scale from 0 to 1. A higher CI value indicates a better fit of the model. equation \ref{equation5} presents the calculation method:
\begin{equation}
\mathrm{CI}=\frac1{Z}\sum_{\delta_i>\delta_j}f(\mathrm{b_i-b_j}),\label{equation5}
\end{equation}
where the affinity value of $\delta_i$ surpasses that of $\delta_j$. Here, $b_i$ and $b_j$ represent the predicted values for $\delta_i$ and $\delta_j$ respectively, while $Z$ denotes the normalization constant. The function $f(\cdot)$ is computed using equation:
$$
f = \begin{cases}
  1 & \text{ if } x>0 \\
  0.5& \text{ if } x=0 \\
  0& \text{ if } x<0
\end{cases}.
$$
\hspace*{1em}The $r_m^2$ metric was proposed in DeepDTA\cite{ozturk2018deepdta} is used to measure the external prediction performance of the model. The metrics are calculated by the formula:
$$
r_m^2=r^2\times(1-\sqrt{r^2-r_0^2})
,$$
where $r$ and $r_0$ are the correlation coefficients with and without intercepts, respectively.

To fairly compare the performance of the models, the splitting methods for the Davis and KIBA datasets are consistent with those used in GraphDTA\cite{nguyen2021graphdta}. The Human and C datasets are split into training, validation, and test sets in a ratio of 8:1:1.
\subsection{Setting of the hyperparameters}
In the HAG-DTA model, the hierarchical GNN network contains two DiffPool layers, which output two coarsened graphs. The first coarsened graph captures the features of functional groups, while the second coarsened graph captures the features of fragment molecules. However, the number of nodes in the coarsened graph needs to be pre-defined, which is an important hyperparameter in the model. For clarity in future discussions, we denote the number of nodes in the first coarsened graph as $n_1$ and the number of nodes in the second coarsened graph as $n_2$. \autoref{table2} defines hyperparameters utilized in the HAG-DTA model. To facilitate comparison of model performance, common hyperparameters for training deep learning models, such as learning rate and number of epochs, are set consistently with GraphDTA\cite{nguyen2021graphdta}. The final three rows in \autoref{table2} represent the search space for the model's hyperparameters.
\begin{table}[!t]
\caption{Hyperparameters used in this study.\label{table2}}%
\begin{tabular*}{\columnwidth}{@{\extracolsep\fill}ll@{\extracolsep\fill}}
\toprule
Hyper-parameters & setting\\
\midrule
Learning rate    & 0.0005\\
Batch size    & 512\\
    $\alpha$       & 0.05  \\
    $\beta$       & 0.05  \\
Dropout rate       & 0.2  \\
\textbf{$n_1$}     & \textbf{4,5,6,7,8,9,10,11,12} \\
\textbf{$n_2$}     & \textbf{2,3}  \\
\textbf{GNN layer} & \textbf{GIN,GCN,GAT,SAGE}\\
\botrule
\end{tabular*}
\begin{tablenotes}%
\item Note: $n_1$ and $n2$ are defined as the number of nodes in the first and second coarsenized graphs, respectively
\end{tablenotes}
\end{table}
\subsection{Hyperparameter Exploration}
The hierarchical GNN module comprises two hyperparameters, namely $n_1$ and $n_2$. In accordance with the design concept of the HAG-DTA model, $n_1$ should be chosen to match the number of functional groups within the drug molecule, while $n_2$ should be consistent with the number of molecular fragments in the drug molecule. However, it is important to note that drug molecules can contain multiple diverse functional groups, and the specific count is contingent on the drug's chemical structure and composition. Even for drugs with an equal number of atoms, the count of functional groups can vary significantly. Additionally, due to the complexity associated with detecting the exact number of functional groups, it is not possible to determine this beforehand for the various drug molecules within the dataset. Therefore, we will explore the impact of different settings of $n_1$ and $n_2$ on the model's performance on four datasets. The experimental setup entails utilizing the GNN layer as GIN and conducting each experiment five times with distinct random seeds. The resulting outcomes are illustrated in \autoref{fig3}.
\newline
\hspace*{1em}For the regression task, the best parameter settings in the Davis dataset are $n_1=4$, $n_2=2$, resulting in an MSE of $0.198 \pm 0.001$ and a CI of $0.908 \pm 0.002$. In the kiba dataset, the best parameter settings are $n_1=6$, $n_2=2$, resulting in an MSE of $0.130 \pm 0.001$ and a CI of $0.892 \pm 0.002$. For the binary classification task, the best parameter settings in the Human dataset are $n_1=7$, $n_2=3$, resulting in an AUC of $0.987 \pm 0.001$ and a Precision of $0.958 \pm 0.005$. In the \textit{C.elegans} dataset, the best parameter settings are $n_1=7$, $n_2=3$, resulting in an AUC of $0.993 \pm 0.001$ and a Precision of $0.963 \pm 0.010$.
\newline
\hspace*{1em}The experimental results from four datasets indicate that the configuration of $n_1$ and $n_2$ sizes significantly impacts the performance and stability of the model. Assigning a value to $n_1$ that exceeds the optimal value leads to instability in the model's performance, resulting in substantial fluctuations. Conversely, setting $n_1$ to a value lower than the optimal value tends to enhance the stability of the model. These findings suggest that excessively large values of $n_1$, surpassing the average number of functional groups in the drug molecules within the dataset, cause instability and prevent accurate capturing of the functional group features. Conversely, excessively small values of $n_1$, falling below the average number of functional groups in the drug molecules within the dataset, although potentially maintaining model stability, lead to poor performance in capturing all the essential features of the functional groups.
\begin{figure*}[!t]
\centering
\includegraphics[width=\linewidth]{image/figure3.pdf}
\caption{Exploration of the effects of $n_1$ (increasing from left to right) and $n_2$ (shown on the legend) settings on experimental results. (a), (b), (c), and (d) represent the experimental results for the Davis, KIBA, Human, and \textit{C.elegans} datasets, respectively.}
\label{fig3}
\end{figure*}
\subsection{Comparisons with other baseline models}
To comprehensively evaluate the performance of our proposed HAG-DTA model, we compare it with the following models: KronRLS and SimBoost based on traditional machine learning methods; TransformerCPI, MolTrans\cite{huang2021moltrans}, MATT\_DTI\cite{zeng2021deep}, DeepDTA, AttentionDTA, and WideDTA based on deep learning methods; GNN-CNN\cite{tsubaki2019compound}, MFR-DTA, TrimNet-CNN, GraphDTA, and IIFDTI\cite{cheng2022iifdti} based on graph-based methods. All the above-mentioned baseline models take sequence features or drug molecule graphs as inputs. Additionally, to ensure a fair comparison of model performances, we exclude DGraphDTA and WGNN-DTA\cite{jiang2022sequence} as they utilize protein-weighted graphs as input for the models.
\begin{table}[t]
\caption{Prediction performance on the Davis and KIBA dataset, sorted by MSE.\label{table3}}
\tabcolsep=0pt%%
\begin{tabular*}{\textwidth}{@{\extracolsep{\fill}}lcccccc@{\extracolsep{\fill}}}
\toprule%
& \multicolumn{3}{@{}c@{}}{Davis} & \multicolumn{3}{@{}c@{}}{KIBA} \\
\cmidrule{2-4}\cmidrule{5-7}%
Model & MSE & CI & $r_m^2$ & MSE & CI & $r_m^2$ \\
\midrule
KronRLS       & 0.379(-)              & 0.871(-)              & 0.407(-)              & 0.411(-)              & 0.782(-)              & 0.342(-)              \\
SimBoost      & 0.282(-)              & 0.872(-)              & 0.644(-)              & 0.222(-)              & 0.836(-)              & 0.629(-)              \\
WideDTA       & 0.262(0.009)          & 0.886(0.003)          & 0.633(0.011)          & 0.179(0.008)          & 0.875(0.001)          & 0.675(0.004)          \\
DeepDTA       & 0.261(0.007)          & 0.878(0.004)          & 0.630(0.017)          & 0.194(0.008)          & 0.863(0.002)          & 0.673(0.019)          \\
MATT\_DTI     & 0.254                 & 0.884(0.004)          & 0.649(0.009)          & 0.151                 & 0.889(0.001)          & 0.745(0.008)          \\
GraphDTA      & 0.229(0.005)          & 0.893(0.002)          & 0.685(0.016)          & 0.139(0.008)          & 0.889(0.001)          & 0.725(0.018)          \\
MFR-DTA       & 0.221(0.001)          & 0.905(0.001)          & 0.705(0.003)          & 0.136(0.001)          & \textbf{0.898(0.002)} & 0.789(0.002)          \\
AttentionDTA  & 0.216(0.019)          & 0.893(0.005)          & \multicolumn{1}{c}{-} & 0.155(0.003)          & 0.882(0.004) & \multicolumn{1}{c}{-} \\
HAG-DTA(GCN)  & 0.201(0.002)          & 0.905(0.001)          & 0.728(0.003)          & \textbf{0.130(0.002)} & 0.888(0.005)          & 0.788(0.006)          \\
HAG-DTA(SAGE) & 0.199(0.003)          & 0.904(0.002)          & 0.730(0.003)          & 0.132(0.002)          & 0.889(0.003)          & \textbf{0.792(0.005)} \\
HAG-DTA(GAT)  & 0.198(0.003)          & 0.907(0.001)          & 0.729(0.002)          & 0.133(0.002)          & 0.890(0.001)          & 0.775(0.004)          \\
HAG-DTA(GIN)  & \textbf{0.198(0.001)} & \textbf{0.908(0.002)} & \textbf{0.733(0.006)} & 0.137(0.004) & 0.886(0.003)          & 0.780(0.003)      \\
\botrule
\end{tabular*}
\begin{tablenotes}%
\item Note: We compare four graph neural network variants: GIN, GAT, GCN, and SAGE. Bold indicates the best result of the evaluation metrics
\end{tablenotes}
\end{table}

\begin{table}[!t]
\caption{Prediction performance on the Human and \textit{C.elegans} dataset, sorted by Precision.\label{table4}}
\tabcolsep=0pt%%
\begin{tabular*}{\textwidth}{@{\extracolsep{\fill}}lcccccc@{\extracolsep{\fill}}}
\toprule%
& \multicolumn{3}{@{}c@{}}{Human} & \multicolumn{3}{@{}c@{}}{\textit{C.elegans}} \\
\cmidrule{2-4}\cmidrule{5-7}%
Model & Precision & Recall & AUC & Precision & Recall & AUC \\
\midrule
TransformerCPI & 0.916(0.006)          & 0.925(0.006)          & 0.973(0.002)          & 0.952(0.006)          & 0.953(0.005)          & 0.988(0.002)          \\
TrimNet-CNN    & 0.918(-)              & 0.953(-)              & 0.970(-)              & 0.946(-)              & 0.945(-)              & 0.987(-)              \\
GNN-CNN        & 0.923(-)              & 0.918(-)              & 0.970(-)              & 0.938(-)              & 0.929(-)              & 0.978(-)              \\
SAG-DTA        & 0.945(0.014)          & 0.933(0.011)          & 0.985(0.002)          & -                     & -                     & -                     \\
IIFDTI         & 0.946(0.017)          & 0.947(0.017)          & 0.984(0.003)          & 0.954(0.010)          & 0.971(0.011)          & 0.991(0.002)          \\
MolTrans       & 0.955(0.012)          & 0.933(0.022)          & 0.974(0.002)          & 0.971(0.007)          & 0.963(0.012)          & 0.982(0.003)          \\
HAG-DTA(GCN)   & 0.942(0.005)          & \textbf{0.952(0.002)} & 0.987(0.002)          & \textbf{0.972(0.009)} & 0.949(0.008)          & 0.990(0.002)          \\
HAG-DTA(SAGE)  & 0.948(0.002)          & 0.945(0.006)          & 0.987(0.001)          & 0.948(0.010)          & 0.950(0.006)          & 0.989(0.001)          \\
HAG-DTA(GIN)   & 0.958(0.005)          & 0.945(0.004)          & 0.987(0.001)          & 0.963(0.010)          & 0.969(0.004)          & \textbf{0.993(0.001)} \\
HAG-DTA(GAT)   & \textbf{0.962(0.004)} & 0.939(0.006)          & \textbf{0.988(0.002)} & 0.948(0.013)          & \textbf{0.973(0.015)} & 0.991(0.002)     \\
\botrule
\end{tabular*}
\begin{tablenotes}%
\item Note: We compare four graph neural network variants: GIN, GAT, GCN, and SAGE. Bold indicates the best result of the evaluation metrics
\end{tablenotes}
\end{table}
According to the results of Section 3.2, we selected the optimal hyperparameter settings and repeated the experiments with five different random seeds to calculate the mean and standard deviation of the results. Meanwhile, for the GNN layer, we experimented with four different GNN frameworks. \autoref{table3} summarizes the model comparison results on the regression dataset, while \autoref{table4} summarizes the model comparison results on the classification dataset. It can be observed that the HAG-DTA method outperforms other baseline models, showing significant improvements in most metrics. Notably, the GIN framework demonstrates superior performance on the HAG-DTA model. However, the differences in model metrics among the four GNN frameworks are inconsequential. These results indicate that the HAG-DTA method proposed in this study captures a more comprehensive set of characteristic information and exhibits superior generalization ability for DTA prediction tasks. The utilization of multiscale feature fusion not only enables precise capture of local structural information related to functional groups in drug molecules but also safeguards the preservation of overall structural details of the compounds.
\subsection{Ablation study}
In our design, the extraction of multi-scale feature information in drug molecules relies on hierarchical GNN networks and deep GNN networks. Layered clustering is employed to extract the local structural information of drug molecules, while the deep GNN networks focus on extracting the global structural information. Furthermore, attention mechanisms are utilized to effectively fuse the multi-scale features. Consequently, to investigate the crucial factors influencing the predictive performance of the model, we generate variants of HAG-DTA through the removal or replacement of specific modules, and subsequently conduct a comprehensive series of ablation experiments in the Davis dataset.
\newline
\hspace*{1em}The results of the ablation experiments are presented in \autoref{table5}. The table clearly demonstrates that the HAG-DTA model exhibits superior performance compared to other variants. Importantly, the performance significantly declines when the attention mechanism or hierarchical GNN network is removed. These findings indicate that incorporating the hierarchical GNN network allows the model to extract more comprehensive topological information from the drug molecules, consequently enhancing the model's predictive capability. Moreover, the attention-based multiscale feature fusion plays a pivotal role in adaptively integrating features, which is essential for affinity prediction.
\begin{table}[!t]
\caption{Performance comparison of ablation experiments.\label{table5}}
\tabcolsep=0pt%%
\begin{tabular*}{\textwidth}{@{\extracolsep{\fill}}lccccccc@{\extracolsep{\fill}}}
\toprule%
& \multicolumn{3}{@{}c@{}}{Module}  & \multicolumn{1}{@{}c@{}}{Loss} & \multicolumn{3}{@{}c@{}}{Result} \\
\cmidrule{2-4}\cmidrule{5-5}\cmidrule{6-8}%
Model & Attention & DGNN & HGNN & $\mathcal{L}_{MMD}$ & MSE & CI & $r_m^2$ \\
\midrule
Model 1      & \checkmark                             & \checkmark                        & ×       & \checkmark                  & 0.208(0.002)          & 0.904(0.003)          & 0.718(0.005)          \\
Model 2      & ×                             & \checkmark                        & \checkmark     & \checkmark                    & 0.207(0.003)          & 0.903(0.002)          & 0.715(0.004)          \\
Model 3      & \checkmark                             & ×    & \checkmark                     & \checkmark                        & 0.203(0.001)          & 0.905(0.002)          & 0.723(0.002)          \\
Model 4 & \checkmark                             & \checkmark                        & \checkmark                       & ×  & 0.201(0.002) & 0.906(0.002) & 0.728(0.004) \\
HAG-DTA(GIN) & \checkmark                             & \checkmark                        & \checkmark     & \checkmark                    & \textbf{0.198(0.001)} & \textbf{0.908(0.002)} & \textbf{0.733(0.006)} \\
\botrule
\end{tabular*}
\begin{tablenotes}%
\item Note: The dataset used for the ablation experiments is Davis. Bold indicates the best result of the evaluation metrics.
\end{tablenotes}
\end{table}

\subsection{Model interpretation}
During the period from 2015 to 2020, analysis of data concerning 164 small molecule drugs approved by the FDA\cite{bhutani2021us} revealed that 87\% (142 drug molecules) possessed N-heterocycles and aromatic skeletons. Chiral molecules accounted for 63\%. In contrast, the occurrence of heterocyclic oxygen compounds was relatively low. These findings suggest a regularity in the composition of drug molecules, where specific structures confer distinct chemical properties to such molecules. Hence, in the realm of computer-aided drug discovery, comprehending how models extract feature information and understanding the impact of these features on the model's performance can offer valuable insights for optimizing the structure of drug molecules.
\newline
\hspace*{1em}The HAG-DTA model utilizes hierarchical GNN networks and deep GNN networks to capture the multi-scale features of drug molecules. By analyzing the cluster allocation matrix and attention scores generated by the model, we can easily understand which functional groups and important atoms are extracted from the drug molecules. We chose the best model trained on the Davis dataset for model interpretation, with hyperparameters $n_1=4$ and $n_2=2$. In this configuration, the first coarsened graph has 4 nodes, while the second coarsened graph has 2 nodes. The model interpretation results for three drug molecules (a), (b), and (c) with CIDs of 9818231, 24779724, and 123631 respectively, in the PubChem database\cite{kim2016pubchem}, are illustrated in Figure 5. Based on the two cluster allocation matrices output by the hierarchical GNN network, we visualized the clustering results. The clustering results of the first layer are labeled with four colors, and the clustering results of the second layer are labeled with black dashed irregular circles, as shown on the left side of \autoref{fig4}. Furthermore, the attention scores for each atom, derived from the output of the deep GNN network, are depicted on the right side of \autoref{fig4}. From the figure, it can be seen that the model extracts the features of the pyrimidine (orange part) and piperidinepropanenitrile (green part) structures of drug molecule (a) during the first layer clustering. During the second layer clustering, the model merges the blue and orange parts and extracts the features of the molecular fragment pyrrole and pyrimidine. Additionally, the model identifies and extracts the functional group information of triazole and quinoline in drug molecule (b), as well as the functional group information of morpholinopropoxy and quinazoline in drug molecule (c). Moreover, the model's global attention scores demonstrate a higher focus on the atoms at the connection points of the two molecular fragments, such as the sulfur atom in drug molecule (b). The above results demonstrate that the HAG-DTA model can accurately identify and extract the features of functional groups and molecular fragments, and it has good interpretability for the DTA task.
\begin{figure}[!t]%
\centering
\includegraphics[width=\linewidth]{image/figure4.pdf}
\caption{The clustering results obtained by HGNN are shown on the left side. The first-layer clustering results are annotated with four different colors, while the second-layer clustering results are annotated with irregular black dashed circles. The attention scores for each atom, derived from the output of DGNN, are depicted on the right side. }\label{fig4}
\end{figure}
\section{Conclusion}
This work presents the HAG-DTA model, which aims to predict Drug-Target Affinity (DTA) by capturing multi-scale feature information of drug molecules. The HAG-DTA model comprises two modules: the hierarchical GNN module and the deep GNN module. The hierarchical GNN module extracts functional groups and molecular fragment features of drug molecules through hierarchical clustering. The deep GNN module employs a gate skipping connection mechanism to fuse multi-level information of molecules and extract global features of drug molecules. Meanwhile, we introduced the MMD loss function to limit the distribution difference between the substructure features and global features. Moreover, the HAG-DTA model incorporates an attention mechanism to adaptively integrate the extracted multi-scale features. Extensive experiments on four publicly available datasets have showcased the superiority of this approach. Furthermore, we performed model interpretation based on hierarchical clustering and attention scores. The results show that the HAG-DTA model exhibits good interpretability, thereby contributing to the advancement of research on drug structure optimization.
\\
\\
\\
\textbf{Abbreviations}
\\
\begin{tabular}{ll} % 左对齐，无边线（简洁风格）
    DTA    & Drug-target affinity \\
    SMILES & Simplified molecular input line entry system \\
    LSTM   & Long short-term memory \\
    GCN    & Graph convolutional networks \\
    GAT    & Graph attention networks \\
    GIN    & Graph isomorphism networks \\
    SAGE    & Graph Sample and Aggregate\\
    CI    & Concordance index \\
    MSE    & Mean squared error \\
\end{tabular}

\section*{Declarations}
\textbf{Ethics approval and consent to participate}
\\
Not applicable.
\\
\\
\textbf{Consent for publication}
\\
Not applicable.
\\
\\
\textbf{Availability of data and materials}
\\
The code and data are provided at \href{https://github.com/peternuo/HAG-DTA}{https://github.com/peternuo/HAG-DTA}.
\\
\\
\textbf{Competing interests}
\\
The authors declare that they have no competing interests.
\\
\\
\textbf{Funding}
\\
This work has received no funding.
\\
\\
\textbf{Author contributions}
\\
J.L. designed and performed the experiments, analyzed the data, and wrote the original manuscript. L.Y. critically reviewed the manuscript, provided substantive edits, and verified the analytical methods. Both authors approved the final version of the manuscript.
\\
\\
\textbf{Acknowledgements}
\\
Not applicable.
\\
\\
\bibliography{sn-bibliography}% common bib file
%% if required, the content of .bbl file can be included here once bbl is generated
%%\input sn-article.bbl

\end{document}
